\documentclass[a4paper,
fontsize=11pt,
%headings=small,
oneside,
numbers=noperiodatend,
parskip=half-,
bibliography=totoc,
final
]{scrartcl}

\usepackage[babel]{csquotes}
\usepackage{synttree}
\usepackage{graphicx}
\setkeys{Gin}{width=.4\textwidth} %default pics size

\graphicspath{{./plots/}}
\usepackage[ngerman]{babel}
\usepackage[T1]{fontenc}
%\usepackage{amsmath}
\usepackage[utf8x]{inputenc}
\usepackage [hyphens]{url}
\usepackage{booktabs} 
\usepackage[left=2.4cm,right=2.4cm,top=2.3cm,bottom=2cm,includeheadfoot]{geometry}
\usepackage[labelformat=empty]{caption} % option 'labelformat=empty]' to surpress adding "Abbildung 1:" or "Figure 1" before each caption / use parameter '\captionsetup{labelformat=empty}' instead to change this for just one caption
\usepackage{eurosym}
\usepackage{multirow}
\usepackage[ngerman]{varioref}
\setcapindent{1em}
\renewcommand{\labelitemi}{--}
\usepackage{paralist}
\usepackage{pdfpages}
\usepackage{lscape}
\usepackage{float}
\usepackage{acronym}
\usepackage{eurosym}
\usepackage{longtable,lscape}
\usepackage{mathpazo}
\usepackage[normalem]{ulem} %emphasize weiterhin kursiv
\usepackage[flushmargin,ragged]{footmisc} % left align footnote
\usepackage{ccicons} 
\setcapindent{0pt} % no indentation in captions
\usepackage{xurl} % Breaks URLs

%%%% fancy LIBREAS URL color 
\usepackage{xcolor}
\definecolor{libreas}{RGB}{112,0,0}

\usepackage{listings}

\urlstyle{same}  % don't use monospace font for urls

\usepackage[fleqn]{amsmath}

%adjust fontsize for part

\usepackage{sectsty}
\partfont{\large}

%Das BibTeX-Zeichen mit \BibTeX setzen:
\def\symbol#1{\char #1\relax}
\def\bsl{{\tt\symbol{'134}}}
\def\BibTeX{{\rm B\kern-.05em{\sc i\kern-.025em b}\kern-.08em
    T\kern-.1667em\lower.7ex\hbox{E}\kern-.125emX}}

\usepackage{fancyhdr}
\fancyhf{}
\pagestyle{fancyplain}
\fancyhead[R]{\thepage}

% make sure bookmarks are created eventough sections are not numbered!
% uncommend if sections are numbered (bookmarks created by default)
\makeatletter
\renewcommand\@seccntformat[1]{}
\makeatother

% typo setup
\clubpenalty = 10000
\widowpenalty = 10000
\displaywidowpenalty = 10000

\usepackage{hyperxmp}
\usepackage[colorlinks, linkcolor=black,citecolor=black, urlcolor=libreas,
breaklinks= true,bookmarks=true,bookmarksopen=true]{hyperref}
\usepackage{breakurl}

%meta
%meta

\fancyhead[L]{U. Wuttke\\ %author
LIBREAS. Library Ideas, 48 (2026). % journal, issue, volume.
\href{https://doi.org/10.18452/37898}{\color{black}https://doi.org/10.18452/37898}
{}} % doi
\fancyhead[R]{\thepage} %page number
\fancyfoot[L] {\ccLogo \ccAttribution\ \href{https://creativecommons.org/licenses/by/4.0/}{\color{black}Creative Commons BY 4.0}}  %licence
\fancyfoot[R] {ISSN: 1860-7950}

\title{\LARGE{Die Genese des Wissenschaftspodcasts \enquote{Die Wissensarchitekt*innen} – von der Idee ins Ohr}}% title
\author{Ulrike Wuttke} % author

\setcounter{page}{1}

\hypersetup{%
      pdftitle={Die Genese des Wissenschaftspodcasts „Die Wissensarchitekt*innen“ – von der Idee ins Ohr},
      pdfauthor={Ulrike Wuttke},
      pdfcopyright={CC BY 4.0 International},
      pdfsubject={LIBREAS. Library Ideas, 48 (2026)},
      pdfkeywords={Wissenschaftskommunikation, Wissenschaftspodcast, Podcasting, Datentracking, science communication, science podcast, data tracking},
      pdflicenseurl={https://creativecommons.org/licenses/by/4.0/},
      pdfcontacturl={http://libreas.eu},
      pdfurl={https://doi.org/10.18452/37898},
      pdfdoi={10.18452/37898},
      pdflang={de},
      pdfmetalang={de}
     }



\date{}
\begin{document}

\maketitle
\thispagestyle{fancyplain} 

%abstracts
\begin{abstract}
\noindent
\textbf{Kurzfassung}: Dieser Beitrag präsentiert in Dialogform als
fiktives Interview die Genese des Podcasts \enquote{Die
Wissensarchitekt*innen}. Diese wird anhand der Punkte Motivation
\& Themenverständnis, Gesprächspartner*innen \& Vorbereitung,
Episodenstruktur \& Dramaturgie, Technik \& Tools und Rezeption,
Feedback \& Ausblick vorgestellt. Dabei werden sowohl spezifische
inhaltliche Aspekte des Podcasts als auch Learnings für allgemein am
Podcasting als Format der Wissenschaftskommunikation interessierte
Personen erörtert.

\begin{center}\rule{0.5\linewidth}{0.5pt}\end{center}
\end{abstract}

%body
\section{Intro}\label{intro}

\textbf{Ulrike Wuttke:} Tapp, tapp, tapp\ldots{} Räusper, ist das
Mikrofon eigentlich an?\footnote{Anmerkung: Dieser Beitrag präsentiert
  in Dialogform als fiktives Interview die Genese des Podcasts
  \enquote{Die Wissensarchitekt*innen}. Die Persona \enquote{Die
  Aufnahmeleitung} ist ein Alter Ego der Autorin und hat keinerlei
  Ähnlichkeiten mit lebenden Personen.}

\textbf{Die Aufnahmeleitung:} \emph{Ja, ich denke schon, wir sind
bereit! Können Sie sich bitte kurz vorstellen?}

\textbf{Ulrike Wuttke:} Ich bin Ulrike Wuttke und ich bin Professorin
für Bibliothekswissenschaft -- Strategien, Serviceentwicklung und
Wissenschaftskommunikation am Fachbereich Informationswissenschaften der
Fachhochschule Potsdam. Februar 2026 bin ich nach langer Vorbereitung
mit einem eigenen Podcast gestartet, er heißt \enquote{Die
Wissensarchitekt*innen}. In ihm rede ich mit spannenden Gäst*innen über
das Thema Datentracking in der Wissenschaft. Danke, dass ich mit dabei
sein darf.

\section{Motivation \&
Themenverständnis}\label{motivation-themenverstuxe4ndnis}

\textbf{Die Aufnahmeleitung:} \emph{Frau Wuttke, Sie haben sich dafür
entschieden, einen neuen bibliothekarischen Podcast zu starten. Auf der
Übersicht vom Bibliotheksverband sind alleine schon über 20
Bibliothekspodcasts verzeichnet}\footnote{O. A. Bibliothekarische
  Podcasts. \emph{Bibliotheksportal}.
  \url{https://bibliotheksportal.de/bibliothekarische-podcasts/}}
\emph{und auch sonst gibt es eine sehr große Anzahl Podcasts. Warum
braucht es noch einen Podcast, nämlich Ihren?}

\textbf{Ulrike Wuttke:} Es kann natürlich nie genügend Podcasts geben!
(lacht). Und auch bibliothekarische Podcasts sind schon eine Weile im
Trend und übernehmen Funktionen sowohl in der Wissenschaftskommunikation
als auch in der Öffentlichkeitsarbeit beziehungsweise im
Marketingmix.\footnote{Siehe zum Beispiel Siegfried, Doreen. 2023.
  Wirkungsorientierte Kommunikationsarbeit: Positionierung der ZBW als
  Open-Science-Partner für die Wirtschaftsforschung. \emph{B.I.T.
  online} 26, Nr. 3: 234--241,
  \url{https://www.b-i-t-online.de/heft/2023-03-fachbeitrag-siegfried.pdf}
  Feidicker, Charlotte. 2025. Ein Medium, das Spaß macht und den
  Austausch fördert SUBTEXT für Partizipation: Podcast \enquote{Über die
  Benutzung von Bibliotheken}. \emph{BuB -- Forum Bibliothek und
  Information} 77, Nr. 10: 496--499,
  \url{https://doi.org/10.24403/jp.1529576}} Daher habe ich letztes Jahr
zusammen mit Alexander Winkler und Ralf Stockmann bei der BiblioCon 2025
in Bremen ein Hands-On-Lab mit dem Titel \enquote{Podcasts in
Bibliotheken} gegeben.\footnote{Wuttke, Ulrike, Ralf Stockmann und
  Alexander Winkler. 2026. Podcasts als bibliothekarisches
  Handlungsfeld. \emph{B.I.T. online} 29, Nr. 1: 48--53,
  \url{https://www.b-i-t-online.de/heft/2026-01-nachrichtenbeitrag-wuttke.pdf}}
Darin ging es neben der eigenen Produktion von Podcasts durch
Bibliothekar*innen auch um bibliothekarische Dienstleistungen rund um
das Podcasting. Mittlerweile hat sogar die ehrwürdige Stabi Berlin im
April 2026 ein eigenes \enquote{Mikrofone ready}-Podcaststudio
eröffnet.\footnote{Bucher, Michael. 2026. Podcaststudio
  Einführungskurse. \emph{Staatsbibliothek zu Berlin: Blog-Netzwerk für
  Forschung und Kultur}. 3. Februar.
  \url{https://blog.sbb.berlin/podcaststudio/}}

\textbf{Die Aufnahmeleitung:} \emph{Und da dachten Sie, warum habe ich
eigentlich keinen Podcast?}

\textbf{Ulrike Wuttke:} Fast! Im Ernst, angesichts der steigenden
Beliebtheit von Podcasts in der Wissenschaftskommunikation habe ich
schon eine Weile mit dem Gedanken gespielt, einen Podcast zu starten, um
mit Expert*innen über Fragen aus dem \enquote{Maschinenraum der
Informationsgesellschaft} zu sprechen, und zwar als Wissenschaftlerin,
aus einer überinstitutionellen Perspektive. Und ich wollte im Rahmen
meiner Transferprofessur mehr Aufmerksamkeit für das Thema Datentracking
in der Praxis erwecken. Da bot es sich an, das Medium Podcast
auszuprobieren, vor allem auch, weil mein Praxispartner für den
Transfer, Ralf Stockmann von der Zentral- und Landesbibliothek Berlin
(ZLB) selbst ein sehr erfahrener Podcaster ist.

\textbf{Die Aufnahmeleitung:} \emph{Und was hat es mit dem Namen des
Podcasts \enquote{Die Wissensarchitekt*innen} auf sich?}

\textbf{Ulrike Wuttke:} Also, ich wollte keinen reinen
Bibliothekspodcast gründen. Ich habe ja auch keine Bibliothek, sondern
einen Podcast, der sich im breiteren Sinn aus Sicht der
Informationswissenschaften mit der Welt hinter den Daten und was
Bibliotheken, Forschung und Algorithmen miteinander verbindet,
beschäftigt. Weil ich selbst ein sehr neugieriger Mensch bin, wollte ich
spannende Personen interviewen, die unsere Informationsinfrastrukturen
gestalten, und mit ihnen kritisch etablierte Strukturen, aktuelle
Entwicklungen und Herausforderungen besprechen. Daher \enquote{Die
Wissensarchitekt*innen}, es sollte auch ein wenig geheimnisvoll klingen.

\textbf{Die Aufnahmeleitung:} \emph{Die erste Staffel Ihres Podcasts ist
dem Thema Datentracking bzw. Science Tracking gewidmet. Warum haben Sie
genau dieses Thema gewählt?}

\textbf{Ulrike Wuttke:} Science Tracking, also Datentracking in der
Wissenschaft insbesondere durch kommerzielle Player, ist ein
hochbrisantes Thema. Was das genau ist, erklären zum Beispiel Renke
Siems und Yuliya Fadeeva in der ersten Episode und Nicola Mößner in der
zweiten Episode, Link ist natürlich in den Shownotes.\footnote{Wuttke,
  Ulrike. 2026. Was ihr schon immer über Datentracking in der
  Wissenschaft wissen wolltet. Mit Renke Siems und Yuliya Fadeeva.
  \emph{Die Wissensarchitekt*innen} (Podcast). 1. Februar 2026.
  \url{https://podcasts.homes/@Die_Wissensarchitekt_innen/episodes/was-ihr-schon-immer-uber-datentracking-in-der-wissenschaft-wissen-wolltet}
  Wuttke, Ulrike. 2026. Business as usual? Philosophische Reflexionen
  zum Datentracking. Mit Nicola Mößner. \emph{Die
  Wissensarchitekt*innen} (Podcast). 1. März 2026.
  \url{https://podcasts.homes/@Die_Wissensarchitekt_innen/episodes/business-as-usual-philosophische-reflexionen-zum-datentracking}}
Angesichts des Ernsts der Lage ist es meine Motivation, zu diesem Thema
einen Transfer in die Gesellschaft und die Praxis herzustellen. Dafür
habe ich von meiner Hochschule, wie bereits erwähnt, eine sogenannte
Transferprofessur bekommen. Hierdurch habe ich Zeit, mich in kooperative
Projekte mit Partner*innen aus Wirtschaft, Politik, Zivilgesellschaft
und Verwaltung einzubringen. Nach langer Recherche und Vorsondierungen,
habe ich mir gedacht, dass so ein Podcast ein großes Potenzial für eine
Breitenwirkung haben kann, um diese Leerstelle im bibliothekarischen
Diskurs zu bespielen, aber auch um darüber hinaus Aufmerksamkeit für
dieses bislang in der Öffentlichkeit unterbelichtete Thema zu
erreichen.\footnote{Weiß, Monika. 2026. Wissenschaft kommuniziert --
  audiovisuell: Science-Videos als digitale Transformation des
  Publizierens. \emph{kommunikation@gesellschaft} 26, Nr. 1.
  \url{https://doi.org/10.15460/kommges.2025.26.1.1866}} Und dann habe
ich gedacht, ok, jetzt probiere ich das aus. So fing eigentlich alles
an.

\textbf{Die Aufnahmeleitung:} \emph{Apropos Transfer und Publikum, an
wen richtet sich der Podcast?}

\textbf{Ulrike Wuttke:} Ich habe ein Bild vor Augen, dass der Podcast
\enquote{Die Wissensarchitekt*innen} ein Muss für alle sein sollte, die
sich für Wissen, Daten und Gesellschaft interessieren und ein Faible für
wissenschaftliche Themen aus Bibliothekswesen, Informationswissenschaft
und Digital Humanities haben. Das können Studierende und
Wissenschaftler*innen in den Informationswissenschaften, Digital
Humanities, Angewandter Informatik und anverwandten Gebieten sein,
Bibliotheksmenschen und Information Professionals, aber auch Personen
aus der Zivilgesellschaft, die sich für Wissen, Daten und Gesellschaft
interessieren. Das Thema ist natürlich schon eine etwas nerdige Nische.
Aber wer sich nun wirklich direkt angesprochen fühlt, muss sich noch
zeigen, der Podcast ist ja gerade erst gestartet. Und gerade für
Nischenthemen eignet sich das Podcastformat durchaus, wie die breite
Themenlandschaft zum Beispiel bei der Plattform
\href{https://wissenschaftspodcasts.de}{Wissenschaftspodcasts.de} zeigt.

\section{Gesprächspartner*innen \&
Vorbereitung}\label{gespruxe4chspartnerinnen-vorbereitung}

\textbf{Die Aufnahmeleitung:} \emph{Wie wählen Sie die
Gesprächspartner*innen aus? Haben Sie da bestimmte Vorgaben oder
Kriterien?}

\textbf{Ulrike Wuttke:} Selbstverständlich müssen die Gäst*innen zum
Thema der Staffel passen. Aber sonst entscheide ich alleine, wen ich
einlade. Die Personen müssen dann natürlich noch \enquote{ja} sagen und
wir müssen einen Termin finden. Ich versuche auch, auf das
Geschlechtergleichgewicht zu achten. Bei technischen Themen treten immer
noch schnell eher männlich gelesene Personen in den Vordergrund. Im
Umkehrschluss heißt es aber nicht, dass ich nur weiblich gelesene
Personen in den Podcast einlade, das ist nicht das Konzept, wie zum
Beispiel bei \enquote{Informatik für die moderne Hausfrau}\footnote{\url{https://informatik-hausfrau.de/}},
ein toller Podcast übrigens.

Ein paar Gäst*innen habe ich schon seit der Planung der ersten Staffel
auf meiner Wunschliste. Momentan wächst diese Liste aber immer weiter,
das liegt auch daran, dass eine feste Rubrik im Podcast die sogenannte
Kettenfrage ist. Da frage ich meine Gäst*innen, wen ich zu welchem Thema
einladen soll. Das ist sehr spannend und inspirierend. Das kann eine
umfangreiche Staffel werden\ldots{}

\textbf{Die Aufnahmeleitung:} \emph{Wie bereiten Sie sich auf die
Gespräche vor? Stellen Sie immer die gleichen Fragen?}

\textbf{Ulrike Wuttke:} Das Gute an einem Podcastinterview ist, dass ich
die Freiheit habe, unterschiedliche und vor allem auch spontane Fragen
zu stellen, auch wenn es feste Rubriken gibt, wie zum Beispiel die eben
erwähnte Kettenfrage. Ich frage auch nicht völlig planlos, sondern
stimme im Vorfeld die Fragen mit den Interviewpartner*innen ab. Das
heißt, ich recherchiere erst umfangreich und führe dann nach Möglichkeit
mit den Gäst*innen ein Vorgespräch zur Feinabstimmung. Gerade Personen
mit wenig Podcasterfahrung wünschen sich das explizit im Vorfeld, das
gibt ihnen eine gewisse Sicherheit. Es hilft auch dabei, einen roten
Faden im Gespräch zu haben. Es ist ja trotzdem immer noch möglich,
nachzuhaken. Gute Podcastinterviews sind also eine Frage der guten
Vorbereitung, nicht von Kontrolle, weil nur dadurch eine gewisse
Spontanität entsteht, die aus meiner Sicht die für ein gutes Gespräch
notwendige Balance zwischen Fachlichkeit und Persönlichkeit schafft.

\textbf{Die Aufnahmeleitung:} \emph{Es gibt also durchaus Raum für
spontane Momente?}

\textbf{Ulrike Wuttke:} Ja, auf jeden Fall! Ich glaube, die Hörer*innen
schätzen an Podcasts als Mittel der Wissenschaftskommunikation, dass sie
weniger \enquote{glatt} produziert sind als typische Radiofeatures.

\section{Episodenstruktur \&
Dramaturgie}\label{episodenstruktur-dramaturgie}

\textbf{Die Aufnahmeleitung:} \emph{Kommen wir noch einmal zur Struktur
des Podcasts. Sie haben ja schon gesagt, dass es feste Bestandteile
gibt. Warum und welche Bestandteile sind das bei den
Wissensarchitekt*innen?}

\textbf{Ulrike Wuttke:} Jede Episode der ersten Staffel hat im Prinzip
eine ähnliche Dramaturgie, um die Hörer*innen mitzunehmen, wobei
Variationen aber durchaus möglich sind. Klassische Bestandteile jedes
Podcasts sind natürlich Intro und Outro, die sozusagen den Rahmen jeder
Episode bilden. Mein Podcast hat ein festes Intro, aber das Outro
spreche ich jedes Mal neu ein, weil ich darin versuche, einen Ausblick
auf die nächste Episode zu geben, sozusagen als Cliffhanger. Nach dem
Intro gibt es die Vorstellung der Gäst*innen, das heißt, erst stelle ich
die Gäst*innen biografisch vor, dann erzählen sie in der Rubrik
\enquote{Persönlicher Pfad} etwas über ihren Hintergrund und wie sie zum
Thema Datentracking gekommen sind, und dann stelle ich spezifische
Fragen zum Fokusthema der Episode, die sich aus Recherche und
Vorgespräch ergeben haben. Dann gibt es meist einen Tooltip, das heißt,
die Gäst*innen stellen ein digitales Werkzeug oder Konzept vor, und die
bereits erwähnte Kettenfrage.

\textbf{Die Aufnahmeleitung:} \emph{Sind alle Fragen vorher
abgesprochen?}

\textbf{Ulrike Wuttke:} Es ist durchaus Platz für Spontanität. Je nach
Konstellation habe ich auch noch die Rubrik \enquote{Gäst*innen fragen}
-- Gäst*innen stellen sich also gegenseitig eine Frage, oder aber auch
mir. Das ist zum Beispiel spontan. Und dann natürlich am Ende, das finde
ich immer sehr höflich, die Frage, ob es noch etwas gibt, was noch nicht
gesagt wurde.

\section{Technik \& Tools}\label{technik-tools}

\textbf{Die Aufnahmeleitung:} \emph{Ich würde gerne noch ein wenig auf
technische und organisatorische Learnings eingehen. Wie läuft die
Produktion so ab, welche Ausstattung ist am Start?}

\textbf{Ulrike Wuttke:} Also vielleicht stelle ich erst einmal klar, der
Podcast \enquote{Die Wissensarchitekt*in\-nen} ist momentan eine
One-Woman-Show und Low-Budget. Ich konzeptioniere, produziere und
veröffentliche alle Episoden alleine und mache Werbung auf Social Media.
Daher ist auch noch nicht alles 100\,\%, aber das Grundprinzip steht.

Zum Schneiden verwende ich die Open-Source-Software
Ultraschall.\footnote{\url{https://ultraschall.fm/}} Das Hosting erfolgt
über \href{https://podcasts.homes}{Podcasts.homes/}.\footnote{\url{https://podcasts.homes/@Die_Wissensarchitekt_innen/about}}
Das ist ein kleiner, aber feiner \enquote{Castopod}-Server\footnote{\enquote{Castopod
  is a free and open-source podcast hosting solution made for podcasters
  who want engage and interact with their audience.} Vgl.
  \url{https://github.com/ad-aures/castopod}} im Fediverse. Ich bin sehr
dankbar, dass der Podcast dort eine Heimat gefunden hat. Ansonsten
spiele ich zusätzliche Materialien, wie zum Beispiel die
Transkriptionen, über meine Webseite aus.\footnote{\url{https://ulrikewuttke.wordpress.com/die-wissensarchitektinnen/}}
Da überlege ich für die Zukunft, ob ich das vielleicht über Zenodo löse,
auch als paralleler Publikationsort für die Mediendateien. Das wäre noch
nachhaltiger und dann bekämen sie zusätzlich einen DOI.

\textbf{Die Aufnahmeleitung:} \emph{Welche Rolle spielt die Technik bei
der Produktion des Podcasts?}

\textbf{Ulrike Wuttke:} Ich habe ein gutes mobiles Aufnahmegerät und ein
Studiomikrofon vom Labor für Wissenschaftskommunikation meines
Fachbereichs zur Verfügung gestellt bekommen, weil eine gute Tonqualität
für ein intimes Medium wie einen Podcast, wo Mensch ja direkt auf den
Ohren der Hörer*innen liegt, sehr wichtig ist. Sonst sehe ich Technik
als Ermöglicherin und nicht als Selbstzweck, wobei mir schon wichtig
ist, weniger proprietäre Tools zu verwenden und datensparsam zu
arbeiten. Da habe ich mich an Good Practices anderer Podcasts orientiert
und viel umgehört, bei Zusammentreffen von Podcaster*innen wie beim
GanzOhr oder einem Workshop, den ich an der FH Potsdam organisiert
habe.\footnote{Ulrike Wuttke (2025), \#WissPodPotsdam2024: Nachbericht
  zum Workshop zu Formaten und Etablierung von Wissenschaftspodcasts
  (Fachhochschule Potsdam, 18.--19.07.2024), DHd-Blog, Link:
  \url{https://dhd-blog.org/?p=21891}} Manchmal sind es so kleine
Sachen, wie darauf zu achten, dass der Podcast einen echten RSS-Feed
hat, damit er plattformunabhängig abonniert werden kann, zum Beispiel
mit der datensparsamen App Antennapod.\footnote{\url{https://antennapod.org/de/}}
Für eine bessere Sichtbarkeit habe ich ihn zusätzlich über die
wichtigsten üblichen Verzeichnisse und Plattformen distribuiert.

\section{Rezeption, Feedback \&
Ausblick}\label{rezeption-feedback-ausblick}

\textbf{Die Aufnahmeleitung:} \emph{Kommen wir nun zu Ihrem Fazit aus
den ersten Episoden. Was hat Sie am meisten überrascht, was gibt es
bislang für Rückmeldungen?}

\textbf{Ulrike Wuttke:} So ein Podcast macht eine Menge Arbeit. Neben
dem Schneiden und den Transkriptionen sind vor allem die Werbung über
Mailinglisten und Social Media zeitintensiv. Positive Rückmeldungen wie
\enquote{Das ist echt unheimlich wichtig.}, \enquote{Erst jetzt habe ich
die Tragweite von Science Tracking verstanden.} oder auch einfach
\enquote{Super interessante Episode!} motivieren aber ungemein und
zeigen meiner Meinung nach, dass das Konzept aufgegangen ist, diesem
nicht nur für Bibliotheksmenschen hochrelevanten Thema einen neuen
Gesprächsraum zu bieten. Vielleicht trägt dazu auch bei, dass ich das
Thema an die Initiative Digital Independence Day angebunden
habe.\footnote{\url{https://di.day/de}}

\textbf{Die Aufnahmeleitung:} \emph{Gibt es noch etwas, das noch gesagt
werden sollte? Welche Pläne haben Sie mit dem Podcast?}

\textbf{Ulrike Wuttke:} Also ein paar weitere Episoden zum Datentracking
habe ich schon in der Planung \ldots{} Ich bin positiv überrascht, wie
viele Personen bereits spontan zugesagt haben. Es besteht natürlich auch
die Möglichkeit, Sonderepisoden zu anderen Themen dazwischenzuschieben
oder weitere Staffeln, auch zu anderen Themen, aufzunehmen. Außerdem
möchte ich gerne die Reichweite erhöhen. Also bitte weitersagen:
\enquote{Die Wissensarchitekt*innen}, Staffel 1 zum Datentracking in der
Wissenschaft, auf \href{https://podcasts.homes}{Podcasts.homes/} und
überall, wo es Podcasts gibt, abonnieren, liken, und eine positive
Bewertung hinterlassen.

Bis dahin, bleibt kritisch -- und vergesst nicht, Wissen ist gestaltbar!

\textbf{Die Aufnahmeleitung:} \emph{Vielen Dank für dieses Interview!}

\textbf{Ulrike Wuttke:} Danke, dass ich hier sein durfte!

\section{Literaturverzeichnis}\label{literaturverzeichnis}

O. A. Bibliothekarische Podcasts. \emph{Bibliotheksportal}.
\url{https://bibliotheksportal.de/bibliothekarische-podcasts/}
(zugegriffen: 14. April 2026).

Bucher, Michael. 2026. Podcaststudio Einführungskurse.
\emph{Staatsbibliothek zu Berlin: Blog-Netzwerk für Forschung und
Kultur}. 3. Februar. \url{https://blog.sbb.berlin/podcaststudio/}
(zugegriffen: 14. April 2026).

Feidicker, Charlotte. 2025. Ein Medium, das Spaß macht und den Austausch
fördert SUBTEXT für Partizipation: Podcast \enquote{Über die Benutzung
von Bibliotheken}. \emph{BuB -- Forum Bibliothek und Information} 77,
Nr. 10: 496--499. \url{https://doi.org/10.24403/jp.1529576}

Siegfried, Doreen. 2023. Wirkungsorientierte Kommunikationsarbeit:
Positionierung der ZBW als Open-Science-Partner für die
Wirtschaftsforschung. \emph{B.I.T. online} 26, Nr. 3: 234--241.
\url{https://www.b-i-t-online.de/heft/2023-03-fachbeitrag-siegfried.pdf}

Weiß, Monika. 2026. Wissenschaft kommuniziert -- audiovisuell:
Science-Videos als digitale \linebreak Transformation des Publizierens.
\emph{kommunikation@gesellschaft} 26, Nr. 1.
doi:\href{https://doi.org/10.15460/kommges.2025.26.1.18666}{10.15460/kommges.20 25.26.1.1866}, \url{https://journals.sub.uni-hamburg.de/hup2/kommges/article/view/1866}
(zugegriffen: 14. April 2026).

Wuttke, Ulrike, Ralf Stockmann und Alexander Winkler. 2026. Podcasts als
bibliothekarisches Handlungsfeld. \emph{B.I.T. online} 29, Nr. 1:
48--53.
\url{https://www.b-i-t-online.de/heft/2026-01-nachrichtenbeitrag-wuttke.pdf}

Wuttke, Ulrike. 2026. Was ihr schon immer über Datentracking in der
Wissenschaft wissen wolltet. Mit Renke Siems und Yuliya Fadeeva.
\emph{Die Wissensarchitekt*innen} (Podcast). 1. Februar 2026.
\url{https://podcasts.homes/@Die_Wissensarchitekt_innen/episodes/was-ihr-schon-immer-uber-datentracking-in-der-wissenschaft-wissen-wolltet}
(zugegriffen: 14. April 2026).

Wuttke, Ulrike. 2026. Business as usual? Philosophische Reflexionen zum
Datentracking. Mit Nicola Mößner. \emph{Die Wissensarchitekt*innen}
(Podcast). 1. März 2026.
\url{https://podcasts.homes/@Die_Wissensarchitekt_innen/episodes/business-as-usual-philosophische-reflexionen-zum-datentracking}
(zugegriffen: 14. April 2026).

Wuttke, Ulrike (2025), \#WissPodPotsdam2024: Nachbericht zum Workshop zu
Formaten und Etablierung von Wissenschaftspodcasts (Fachhochschule
Potsdam, 18.--19.07.2024), \emph{DHd-Blog}, Link:
\url{https://dhd-blog.org/?p=21891} (zugegriffen: 14. April 2026).

%autor
\begin{center}\rule{0.5\linewidth}{0.5pt}\end{center}

\textbf{Ulrike Wuttke} (\url{https://orcid.org/0000-0002-8217-4025}) ist
Professorin für Bibliothekswissenschaft, Strategien, Serviceentwicklung
und Wissenschaftskommunikation an der Fachhochschule Potsdam. Seit April
2025 hat sie an der FH Potsdam eine Transferprofessur zum Thema
Datentracking.

\end{document}